\begin{figure*}[t]
\centering
\includegraphics[width=\textwidth]{figures/v33_structural_ablation.pdf}
\caption{Component ablation on the same 354 targeted graph mutants. (a) Escapes divided by the target set of each removed module; different target counts do not rank module importance. (b) Thirteen mutation types against eight validator variants, with accepted/targeted counts in each cell. C1--C6 follow the module order in (a); None removes all six. The remaining twelve missing-dwell cases are protected by semantic checks. These are constructed configuration faults, not additional language items or flights.}
\label{fig:v33-structural-ablation}
\end{figure*}

\begin{figure*}[t]
\centering
\includegraphics[width=\textwidth]{figures/v33_fair_language_ablation.pdf}
\caption{Equal-vocabulary, anonymous-ID information ablation on nineteen unchanged texts, with five isolated contexts per condition. Small points in (a) are individual runs; large points are means. The reviewed reference contains five abstract height upper-bound items and fourteen unsupported items; repeated contexts do not increase the nineteen-item corpus. (b) Raw and structurally admitted false acceptances use 70 out-of-scope opportunities per condition. (c) Reference-positive gate coverage uses 25 opportunities; it is descriptor acceptance, not executable physical-monitor coverage. Without skeletons, strict canonical serialization blocks every descriptor while tuple selection remains high.}
\label{fig:v33-fair-language-ablation}
\end{figure*}

% Supplement only:
\begin{figure}[htbp]
\centering
\includegraphics[width=\textwidth]{figures/v33_binding_ambiguity.pdf}
\caption{Candidate assignments for the twenty-four frozen internal formula records. The logarithmic axis compares explicit role bindings with coarse adjacency by kind and unit. The latter jointly removes several binding safeguards; this is a configuration-ambiguity stress test rather than an isolated causal test of one graph feature.}
\label{fig:v33-binding-ambiguity}
\end{figure}
